\chapter{Marco teorico y antecedentes}

\section{Accesibilidad y comunicacion en Lengua de Señas Argentina}

Completar con una descripcion de la LSA, su importancia como lengua natural de la comunidad sorda argentina y las barreras comunicacionales que pueden aparecer en la interaccion con personas oyentes.

\section{Sistemas de traduccion de señas}

Completar con antecedentes sobre sistemas de reconocimiento y traduccion de lenguas de señas, diferenciando enfoques basados en video, sensores, guantes, landmarks, glosas o representaciones intermedias.

\section{Vision por computadora para reconocimiento gestual}

Completar con los conceptos principales de vision por computadora aplicados al reconocimiento de gestos: captura de imagen, deteccion de manos, seguimiento temporal y representacion de movimiento.

\section{Landmarks de manos}

Completar con la idea de landmarks como puntos de referencia que representan la posicion de articulaciones de la mano. En este proyecto, estos puntos permiten transformar una imagen en una representacion numerica apta para un modelo de clasificacion.

\section{Modelos livianos en dispositivos moviles}

Completar con la importancia de ejecutar modelos livianos en dispositivos moviles, considerando restricciones de CPU, memoria, bateria, latencia y disponibilidad offline.

\section{Limitaciones de los enfoques existentes}

Completar con las principales limitaciones identificadas: escasez de datasets en LSA, dependencia de servidores externos, problemas de privacidad, necesidad de conectividad y dificultad de generalizar gestos entre distintas personas.
